% Ч1.5 — Перевод единиц измерения (скорость, время, длина, площадь, объём)
% Старая методичка: стр. 80; 92 (старые задания 9; 10)
% Источники: методичка по заданию 11 (движение). Стиль — стиль-изложения.md

\begin{tema}{Ч1.5}
\opr{Скорость.} $1\ \text{м/с}=\dfrac{0{,}001\ \text{км}}{\frac{1}{3600}\ \text{ч}}=3{,}6\ \text{км/ч}$.
Поэтому:
\begin{itemize}
  \item чтобы перевести м/с в км/ч, умножьте на 3,6;
  \item чтобы перевести км/ч в м/с, разделите на 3,6.
\end{itemize}

\opr{Время.} Минуты переводите в часы \emph{обыкновенной дробью}:
$20\ \text{мин}=\frac{20}{60}\ \text{ч}=\frac13\ \text{ч}$.
\[
  15\ \text{мин}=\tfrac14\ \text{ч},\quad 20\ \text{мин}=\tfrac13\ \text{ч},\quad
  40\ \text{мин}=\tfrac23\ \text{ч},\quad 45\ \text{мин}=\tfrac34\ \text{ч}.
\]

\begin{pfTablica}{lll}
Длина & Площадь & Объём\\
$1\ \text{км}=1000\ \text{м}$ & $1\ \text{м}^2=10\,000\ \text{см}^2$ & $1\ \text{м}^3=1000\ \text{л}$\\
$1\ \text{м}=100\ \text{см}$ & $1\ \text{км}^2=10^6\ \text{м}^2$ & $1\ \text{л}=1\ \text{дм}^3=1000\ \text{см}^3$\\
$1\ \text{см}=10\ \text{мм}$ & $1\ \text{га}=10\,000\ \text{м}^2$ & $1\ \text{мл}=1\ \text{см}^3$\\
\end{pfTablica}

\begin{primer}
Переведите 72 км/ч в м/с. Решение: $72:3{,}6=20$. Ответ: 20 м/с.
\end{primer}

\begin{oshibka}
$20\ \text{мин}\ne0{,}2\ \text{ч}$. Правильно: $20\ \text{мин}=\frac13\ \text{ч}$.
Десятичная запись $0{,}33\ldots$ неточна — оставляйте дробь.
\end{oshibka}
\end{tema}
